% =============================================================================================
% §2 Related work. Four clusters, each closing on what this paper does differently.
%
% Every entry cited here had its metadata verified against the arXiv abstract page or the
% publisher record on 2026-08-10; refs.bib records which venue fields were confirmed and which
% were left in preprint form rather than guessed.
% =============================================================================================
\section{Related work}
\label{sec:related}

\subsection{Foundation models for time series, and for K-lines}
\label{sec:rel-foundation}

A line of recent work trains a single sequence model on large, heterogeneous collections of time
series and evaluates it zero-shot or with light adaptation. \citet{ansari2024chronos} discretize
scaled real values into a fixed token vocabulary and train a language model over the resulting
symbols; \citet{das2024timesfm} train a decoder-only model over patched inputs;
\citet{woo2024moirai} train across frequencies and domains with a unified objective. Upstream of
these, \citet{nie2023patchtst} established the patch-as-token treatment that much of the line
inherits. The shared premise is that a general representation of temporal structure transfers.

\citet{shi2025kronos} apply that premise to financial K-lines, and it is the design this paper
starts from: a tokenizer trained to reconstruct per-bar market features, followed by an
autoregressive model over the resulting discrete codes. We inherit the two-stage arrangement, the
backbone dimensions and much of the training and evaluation scaffolding, and we vary one thing.
Every component here is written from scratch; the parent's weights are not used, and no result in
this paper is a comparison against them.

We should be exact about what this artifact is, because the surrounding literature invites a
larger claim than we can support. This is not a foundation model. The backbone is a fixed,
$31{,}795{,}200$-parameter measurement instrument, held identical across every arm so that the
quantizer and input arm are the only varied factors. It is trained on one venue, one instrument
class and one frequency, and we make no transfer claim of any kind. The object under study is the
tokenizer.

\subsection{Discrete tokenization, and what its objective actually optimizes}
\label{sec:rel-tokenization}

Learned discrete representations descend from \citet{oord2017vqvae}, which pairs an encoder with a
learned codebook and a commitment loss. Two later families remove the codebook.
\citet{mentzer2024fsq} quantize each latent dimension to a small fixed set of levels, which
eliminates the codebook, the commitment term and codebook collapse together, at the cost of a
vocabulary that is a product of per-dimension level counts. \citet{zhao2024bsq} project to a
hypersphere and binarize, giving an implicit vocabulary exponential in dimension.
\citet{yu2024lfq} report that tokenizer design, rather than the generative model above it, is what
determines downstream quality --- a finding adjacent to ours in spirit and opposite in method,
since it is established by improving a tokenizer rather than by measuring what one discards.

The relevant observation for this paper is about evaluation rather than architecture. These
methods are compared on reconstruction fidelity, codebook utilization and downstream generation
quality, and the first of those is the objective they are trained on. That is appropriate where
reconstruction is the goal --- an image tokenizer that reconstructs well has done its job. It
becomes an assumption when the codes are handed to a forecaster, because the quantity that matters
downstream is not how much of the input survives but \emph{which parts} do. Reconstruction error
is an average over dimensions weighted by their contribution to that average, so a dimension can
be discarded almost entirely while the aggregate metric barely moves.

We are not aware of work that measures, for a reconstruction-trained tokenizer, which input
dimensions are preferentially discarded and whether the discarded ones are the task-relevant ones.
\Cref{sec:mechanism} is that measurement, on a fixture built so the answer is checkable, and
\cref{sec:mech-fix} is an interface that changes it.

\subsection{Microstructure as a predictive signal, and the distinction we are careful about}
\label{sec:rel-micro}

Order-flow imbalance is among the most robust short-horizon predictors in the microstructure
literature. \citet{cont2014ofi} show that price changes over short intervals are driven largely by
the imbalance between order-book events at the best quotes, a linear relationship with substantial
explanatory power. Related work reads informational content from the composition of flow rather
than its direction: \citet{easley2012vpin} construct a toxicity measure from volume-bucketed trade
imbalance and relate it to liquidity provision.

The distinction that governs this paper is that these are not one quantity. Order-flow imbalance
in the sense of \citet{cont2014ofi} is computed from \emph{order-book depth} --- the arrival,
cancellation and execution of resting liquidity at the best quotes. What we compute is signed
\emph{executed-volume} imbalance from aggregated trades, which requires no depth feed and is
available in the public dumps this study is built on. Trade signing itself has a literature
\citep{lee1991trade}; our source labels the aggressor side directly, so no signing heuristic is
applied and none of that literature's classification error enters our features.
% artifact: src/trikaal/data/reduce.py:75 — buy = pl.col("is_buyer_maker").not_(), i.e. taker BUY is
%   read from the exchange's own flag on each aggregated trade, not inferred from a quote rule.

We therefore use the name trade-flow imbalance throughout, in prose, configuration and code, and
never the order-flow name. This is not pedantry about terminology. The published effect sizes
belong to the depth-based quantity, they are not evidence about ours, and a negative result here
is evidence about trade-flow imbalance only. Order-book depth is explicit future work
(\cref{sec:lim-features}).

\subsection{Evaluation discipline}
\label{sec:rel-eval}

The methodological literature this paper leans on is concerned with a single failure: that a
sufficiently determined search over strategies will produce an impressive backtest from noise.
\citet{harvey2016crosssection} argue that the accumulated multiplicity of published cross-sectional
predictors demands a substantially raised significance bar, and that most published findings would
not survive one. \citet{bailey2014dsr} give the deflated Sharpe ratio, which discounts an observed
Sharpe by the number of trials it was selected from and by the non-normality of the return series.
\citet{lopezdeprado2018advances} supplies the sample-hygiene machinery --- purging, embargoing and
combinatorially purged cross-validation --- that keeps overlapping labels from leaking a training
window into its own evaluation.

Our design takes all four instruments: purged walk-forward with a measured embargo, a paired
moving-block bootstrap, a deflated Sharpe ratio against an enumerated trial count, and a
cost-aware net information ratio rather than a gross one. \Cref{sec:setup} specifies them.

What we add is not a new statistic but a procedural commitment. The deflated Sharpe ratio corrects
for a trial count, and the trial count is supplied by the researcher; nothing in the instrument
prevents an author from reporting a smaller one than they searched. Pre-registration is what makes
that number checkable, and it is rare in this literature. We fixed the decision rule, the
thresholds, the horizon, the cost, the seed count and the trial enumeration before any of their
inputs existed, published the amendment log with the direction of each change recorded before the
change was evaluated, and pre-committed a null as publishable. \Cref{sec:repro} describes the
machinery; \cref{sec:repro-defects} reports what it caught.

\subsection{Position}
\label{sec:rel-position}

Against that background this paper makes one claim. Reconstruction-trained tokenizers allocate
code capacity by variance and covariance and never by downstream value, so a task-relevant channel
that is low-variance and weakly covariant is excluded --- deterministically, under conditions
\cref{sec:mech-claim} states --- and microstructure is exactly such a channel. The remedy is an
interface property rather than a bigger model or a longer schedule. Whether the same effect governs
real trade-flow imbalance in a real market is a separate question, and it is the one
\cref{sec:setup} was pre-registered to answer.
